%% ==========================================================================
%% DRAFT -- August 12, 2026. Author byline: Daniel Kirtchakov (05oz,
%% halfounce.io; no institutional affiliation; ORCID 0009-0009-5213-4098).
%%
%% Framing constraints honoured (PROTOCOL section 11): the note opens with the
%% mathematical question, states every certified object precisely, separates
%% what is certified from what is enclosed and what is trusted, reports the
%% pre-registered kill condition and its measured statistics as plainly as the
%% positive results, cites the computation in a single paragraph, and closes
%% with the questions it opens.
%% ==========================================================================
\documentclass[11pt]{amsart}

\usepackage[margin=1in]{geometry}
\usepackage{amsmath,amssymb}
\usepackage{booktabs}
\usepackage{array}
\usepackage[colorlinks=true,allcolors=blue]{hyperref}

\emergencystretch=3em
\tolerance=1200
\hbadness=4000

\newcommand{\mT}{\mathrm{mT}}
\newcommand{\MHz}{\mathrm{MHz}}
\DeclareMathOperator{\diag}{diag}
\DeclareMathOperator{\dist}{dist}
\DeclareMathOperator{\tr}{tr}

\theoremstyle{plain}
\newtheorem{theorem}{Theorem}[section]
\newtheorem{lemma}[theorem]{Lemma}
\newtheorem{proposition}[theorem]{Proposition}
\newtheorem{corollary}[theorem]{Corollary}
\theoremstyle{definition}
\newtheorem{definition}[theorem]{Definition}
\theoremstyle{remark}
\newtheorem{remark}[theorem]{Remark}
\newtheorem{question}[theorem]{Question}

\begin{document}

\title[Certified brackets at the ZEFOZ points of
$^{167}$Er$^{3+}$:Y$_2$SiO$_5$]{Certified brackets at the published ZEFOZ
points of $^{167}$Er$^{3+}$:Y$_2$SiO$_5$: existence, curvature, and the
measured cost of a completeness certificate}

\author{Daniel Kirtchakov}
\address{Independent researcher (\texttt{05oz}), Half Ounce Research; no
institutional affiliation. ORCID 0009-0009-5213-4098.
\texttt{daniel@halfounce.io} $\cdot$ \texttt{https://halfounce.io}}

\date{Draft of August 12, 2026.}

\begin{abstract}
The magnetic-field points at which a hyperfine transition frequency of
$^{167}$Er$^{3+}$:Y$_2$SiO$_5$ is stationary --- the ZEFOZ
(zero-first-order-Zeeman) points --- set the predicted coherence times of a
leading solid-state quantum-memory platform.  The published atlas of these
points (Matsuura et al., Phys.\ Rev.\ B \textbf{113}, 085421 (2026);
arXiv:2412.10126) is produced by Newton iteration from finite grids of starting
fields, and its authors state plainly that the number of points found depends
on the initial grid.  Two questions follow.  First: do the published points
survive rigorous scrutiny --- does a true stationary point of the exact spin
Hamiltonian provably exist near each tabulated field, and can its location,
frequency, and curvature (the quantity that sets the predicted $T_2$) be pinned
by two-sided bounds a skeptic can re-derive with a standard-library program?
Second: can the published list be proven \emph{complete} over an explicit field
box at laptop scale?  To the first question we give a full positive answer: for
each of the twenty tabulated nonzero-field points of both crystallographic
sites we certify, in exact rational arithmetic, eigenvalue and transition
brackets of width $4\times10^{-10}$~MHz, a gradient-norm bound
$|\nabla f|\le 3.2\times10^{-37}$~MHz/mT at an exactly specified rational
field, two-sided brackets on all three eigenvalues of the $3\times3$ frequency
Hessian, its signature, and --- by a Krawczyk contraction with a rigorous
third-derivative bound --- the existence and local uniqueness of an exact
stationary point within $2.9\times10^{-14}$~mT of the stated field.  The
certified signatures show that none of the twenty published points is a local
minimum of its transition frequency: thirteen are saddles and seven are local
maxima.  At zero field we certify exactly, by an explicit time-reversal
identity, that all 120 transitions of both sites are stationary, with certified
curvature brackets for the ten published zero-field pairs.  To the second
question the answer is a documented negative: a pre-registered kill condition
for the completeness search fired.  Certified branch-and-bound exclusion over
the box $\|B\|_\infty\le100$~mT, measured on stratified sample chunks, gives a
strict lower bound above $10^{6}$ laptop-hours against a pre-registered budget
of $200$; the
obstruction --- quasi-degenerate hyperfine doublets that force per-level
spectral-gap machinery below its validity radius --- is quantified.  Three
errata in the reference paper are recorded, with version history: a
load-bearing sign error in a printed quadrupole matrix that stood for eleven
months (v1, v2; corrected upstream in v3), two tabulated frequencies
inconsistent with their own stationary points by $2.7$ and $4.3$~MHz
(certified restatements given), and an inconsistent sign pairing in a printed
field vector.  The certificate and a standard-library checker are the public
unit.
\end{abstract}

\maketitle

\section{The question}\label{sec:question}

A hyperfine transition of an $^{167}$Er$^{3+}$ ion in Y$_2$SiO$_5$ has a
frequency $f_{ij}(B)=\lambda_j(B)-\lambda_i(B)$, the difference of two
eigenvalues of a $16\times16$ effective spin Hamiltonian depending on the
applied magnetic field $B\in\mathbb R^3$.  At a field where $\nabla
f_{ij}(B)=0$ --- a ZEFOZ point \cite{Fraval2004} --- the transition is
first-order insensitive to magnetic noise, and the residual dephasing is set by
the second-order response, i.e.\ by the $3\times3$ Hessian of $f_{ij}$.
Operating at such points is the standard route to long spin coherence in
rare-earth quantum memories \cite{Zhong2015,Ortu2601}, and for
$^{167}$Er$^{3+}$:Y$_2$SiO$_5$ --- a rare-earth system combining a
telecom-band optical interface with a nuclear spin --- the atlas of ZEFOZ
points and their curvatures published by Matsuura, Yasui, Kaji, Sasakura,
Tawara, and Adachi \cite{Matsuura} is the design input for spin-wave storage
proposals \cite{Matsuura,Business2606}.

That atlas is numerical in an essential way.  The points are found by Newton
iteration from finite grids of initial fields, and the authors state the
consequence themselves: ``because the total number of ZEFOZ points found
depends on the initial grid conditions, the initial magnetic fields for the
search were established by combining the three conditions below''
\cite[App.~D]{Matsuura}.  No error bounds accompany the tabulated locations,
frequencies, or curvatures; nothing excludes further points; and, as we found
in the course of this work (Section~\ref{sec:errata}), the printed inputs of
the calculation carried a load-bearing sign error through two arXiv versions.

We therefore ask two questions, of different difficulty.

\begin{question}\label{q:certify}
Near each published ZEFOZ field, does a stationary point of the \emph{exact}
transition frequency provably exist?  Can its location, its frequency, and the
full spectrum of its Hessian be enclosed in two-sided brackets that are
machine-checkable from first principles --- exact rational arithmetic on the
published Hamiltonian --- by a program short enough to audit, with no trust in
the software that produced the brackets?
\end{question}

\begin{question}\label{q:complete}
Can the published list be proven \emph{complete} --- every stationary point of
every one of the 120 transitions enumerated --- over an explicit field box, at
laptop scale?
\end{question}

Question~\ref{q:certify} we answer in full (Theorems~\ref{thm:zero} and
\ref{thm:points}).  Question~\ref{q:complete} carried a pre-registered kill
condition, and the kill condition fired; Section~\ref{sec:kill} reports the
measured exclusion statistics that fired it and the structural obstruction they
expose.  Both outcomes are results, and both are recorded with the same
precision.

\section{The model and its exact rationalization}\label{sec:model}

The ground-manifold ($^{4}I_{15/2}$, $Z_1$) effective Hamiltonian of
$^{167}$Er$^{3+}$:Y$_2$SiO$_5$ used throughout \cite{Matsuura} is
\begin{equation}\label{eq:H}
H(B) \;=\; \mathbf I\cdot A\,\mathbf S \;+\; \mathbf I\cdot Q\,\mathbf I
\;+\; \mu_B\, B\cdot g\,\mathbf S \;-\; \mu_N\, g_n\, B\cdot \mathbf I ,
\end{equation}
with electron spin $S=1/2$, nuclear spin $I=7/2$ (dimension $16$), field $B$ in
the optical frame $(D_1,D_2,b)$, and $A$, $Q$, $g$ the site-dependent
$3\times3$ matrices printed in \cite[App.~B]{Matsuura} (v3), attributed there
to Wang et al.\ \cite{Wang2023}.  We adopt every printed entry exactly as a
rational number, take $\mu_B/h=13.996244936$~MHz/mT and
$\mu_N/h=7.622593285\times10^{-3}$~MHz/mT (CODATA truncations, again as exact
rationals), $g_n=-0.1618$ as printed, and both crystallographic sites.  The
only irrationalities in $H(B)$ are the nuclear ladder amplitudes $\sqrt7$,
$\sqrt{12}$, $\sqrt{15}$; the certificate encloses them in rational intervals
validated by squaring, and all further arithmetic is exact rational
rectangle-interval arithmetic with outward rounding.  Levels
$\lambda_0\le\dots\le\lambda_{15}$ are indexed ascending, $f_{ij}=\lambda_j-
\lambda_i$, and all bracket claims below are claims about the exact spectrum of
the exact matrix \eqref{eq:H}: every stated interval contains the corresponding
true value.  (Whether \eqref{eq:H} with these matrices describes the physical
crystal is a separate, experimental question; the certificate is unconditional
about the mathematics and silent about the spectroscopy.)

\section{Certified statements}\label{sec:results}

\subsection{Zero field}\label{sec:zero}

\begin{theorem}[exact zero-field stationarity]\label{thm:zero}
Let $M\in\{0,\pm1\}^{16\times16}$ be the signed permutation
$M=R^{(1/2)}\otimes R^{(7/2)}$, $R^{(j)}\lvert j,m\rangle=(-1)^{\,j-m}\lvert
j,-m\rangle$.  Then, as an exact identity of matrices over
$\mathbb Q[\sqrt7,\sqrt{12},\sqrt{15}]$, for both sites,
\[
M\,\overline{H_0}\,M^{\mathsf T}=H_0,\qquad
M\,\overline{Z_k}\,M^{\mathsf T}=-Z_k\quad(k=1,2,3),
\]
where $H_0=H(0)$ and $Z_k=\partial H/\partial B_k$.  Consequently
$H(-B)=MH(B)^{\mathsf T}M^{\mathsf T}$, the spectrum of $H(B)$ is even in $B$;
and since the certificate also proves the sixteen zero-field levels of each
site simple (disjoint brackets of width $2\times10^{-10}$~MHz), every level is
analytic near $B=0$ and
\[
\nabla f_{ij}(0)=0 \qquad\text{exactly, for all } 0\le i<j\le15,
\text{ both sites.}
\]
\end{theorem}

This upgrades the folklore statement that ``at zero field all transitions are
ZEFOZ'' to a machine-checked identity: the checker verifies the symbolic
transform entry-by-entry in $\mathbb Q[\sqrt d]$ and the simplicity via
Sylvester inertia counts (below).  For the ten zero-field pairs tabulated in
\cite[Table~4]{Matsuura} the certificate additionally carries two-sided
brackets on all three eigenvalues of each Hessian $\nabla^2 f_{ij}(0)$
(Section~\ref{sec:method}); for example, for site~1 $(i,j)=(7,9)$ the certified
Hessian spectrum is $\{-11.9905\ldots,\,+1.6899\ldots\times10^{-4},\,
+1.1300\ldots\}$~MHz/mT$^2$, each bracket of width below $2.1\times10^{-9}$.  The
signatures are mixed: six of the ten published zero-field transitions have an
indefinite Hessian, so at zero field those six frequencies are saddles in $B$;
the remaining four --- site~1 $(6,8)$ and $(6,9)$, site~2 $(6,8)$ and $(7,11)$
--- carry the certified signature $(+,+,+)$, so their Hessians are
positive-definite and at zero field those four frequencies are local minima
in $B$.

\subsection{The twenty published nonzero-field points}\label{sec:points}

\begin{theorem}[certified brackets, curvature, and existence]\label{thm:points}
For each of the twenty entries of \cite[Table~5]{Matsuura} (ten per site)
there is an exactly specified rational field $B^*$ (dyadic, denominator
$2^{120}$), within the rounding radius of the published coordinates, such that
the certificate proves:
\begin{enumerate}
\item[(i)] brackets of width $2\times10^{-10}$~MHz for all sixteen
$\lambda_n(B^*)$, hence a bracket of width $4\times10^{-10}$~MHz for
$f_{ij}(B^*)$;
\item[(ii)] $\lvert\nabla f_{ij}(B^*)\rvert \le 3.2\times10^{-37}$~MHz/mT;
\item[(iii)] two-sided brackets, of width at most $2.1\times10^{-15}$~MHz/mT$^2$,
for the three eigenvalues of $\nabla^2 f_{ij}(B^*)$, with certified signs;
\item[(iv)] existence and uniqueness of a point $\hat B$ with $\nabla
f_{ij}(\hat B)=0$ in the box $B^*+[-r,r]^3$, $r=2^{-45}\,\mT\approx
2.9\times10^{-14}\,\mT$ (Krawczyk contraction; worst-case contraction ratio
$0.096$ over the twenty points).
\end{enumerate}
\end{theorem}

\begin{corollary}[certified stationary-point types]\label{cor:sig}
Of the twenty published ZEFOZ points, none is a local minimum of its
transition frequency.  All ten site-1 points and the three site-2 points
$(5,6)$, $(5,7)$, $(6,7)$ have certified signature $(-,-,+)$ (saddles); the
remaining seven site-2 points, including the optimal $(14,15)$ transition of
\cite{Matsuura}, have signature $(-,-,-)$ (local maxima).
\end{corollary}

The Hessian brackets are the certified form of the quantity the application
consumes: in the noise model of \cite[Eq.~(8)]{Matsuura} the predicted $T_2$ at
a ZEFOZ point is a function of the second-order response along the noise
distribution, i.e.\ of the enclosed Hessian; every curvature number a
proposal takes from \cite{Matsuura} can now be taken with a proof.  For the
site-2 $(14,15)$ point singled out in \cite{Matsuura} the certified Hessian
spectrum is
$\{-1.65241\times10^{-4},\,-8.94878\times10^{-5},\,-1.44676\times10^{-6}\}$
MHz/mT$^2$ (brackets of width $<3\times10^{-16}$), and the certified field is
$\hat B=(-378.9856,\,+73.2675,\,+502.3521)$~mT $\pm\,2.9\times10^{-14}$~mT ---
which also settles the sign question of Section~\ref{sec:errata}, item (E3).

\subsection{Method, in brief}\label{sec:method}

All claims reduce to rational inequalities verified by the checker; the
mathematical ingredients are classical.  Eigenvalue brackets come from
interval $LDL^{\mathsf T}$ factorizations of $H(B^*)-\mu I$ at certified
shifts: when every pivot interval is sign-definite, Sylvester's law of inertia
proves the number of eigenvalues below $\mu$.  High-precision eigenvectors
(computed to $70$ digits, then frozen as dyadic rationals with denominator
$2^{170}$ --- candidates, not trusted) are certified a posteriori: the exact
interval residual $\|H\hat v-\tilde\lambda\hat v\|$ against the certified gaps
gives Davis--Kahan $\sin\theta$ bounds \cite{DavisKahan} of order $10^{-42}$;
Hellmann--Feynman then turns them into gradient enclosures (ii), and
second-order perturbation sums with certified matrix elements and certified
spectral gaps give a signed enclosure of the full Hessian (iii), whose
eigenvalues are bracketed by Gershgorin discs after an approximate rotation,
with the exact Ostrowski congruence correction \cite[Thm.~4.5.9]{HornJohnson}
for the rotation's dyadic non-orthogonality.  Existence (iv) is the Krawczyk
test \cite{Krawczyk,Neumaier}: with $C\approx(\nabla^2f)^{-1}$ rational and
$X=B^*+[-r,r]^3$,
\[
K(X)=B^*-C\,\nabla f(B^*)+\bigl(I-C\,\nabla^2 f(X)\bigr)\,[-r,r]^3
\subset\operatorname{int}X
\]
proves a unique zero of $\nabla f$ in $X$.  The Hessian enclosure over $X$ is
the center enclosure inflated entrywise by $9\sqrt3\,Mr$, where
$M=192\,D^3(\gamma_i^{-2}+\gamma_j^{-2})$ is a rigorous bound on all third
directional derivatives of $f_{ij}$ over $X$, derived in
Appendix~\ref{app:third} from a resolvent estimate and a Cauchy coefficient
bound; $D$ bounds $\sup_{|e|=1}\|\partial_e H\|_2$ and $\gamma_{i},\gamma_j$
are certified gap lower bounds over $X$.  The margins are enormous ---
$\lvert\nabla f(B^*)\rvert\lesssim10^{-37}$ against a needed
$\lesssim10^{-18}$ --- because the flatness of clock transitions makes the
Krawczyk test brutally stringent: the smallest certified Hessian eigenvalue
magnitude among the twenty points is $5.67\times10^{-8}$~MHz/mT$^2$, so the
preconditioner norm reaches $1.8\times10^{7}$ and existence at these points is
provable only from a candidate polished far beyond double precision.  A
Newton--Kantorovich argument with double-precision data alone would fail by
many orders of magnitude; this appears to be intrinsic to certifying clock
points, not an artifact of our route.

\subsection{Relation to the pilot certificate and to prior tools}

A same-program pilot release certified eigenvalue and transition brackets at
these fields (without stationarity, curvature, or existence); its certificate
re-verified here (checker exit 0; 352 brackets against an independent 60-digit
diagonalization, zero containment failures).  The present certificate strictly
extends it.  Generic verified eigensolvers exist --- symveig \cite{symveig}
encloses spectra of fixed Hermitian matrices, and verified-numerics methodology
is standard \cite{Rump} --- but we find no prior certified treatment of
\emph{parametric} eigenvalue-difference stationary points (gradients, Hessians,
existence, or completeness) for spin Hamiltonians, and no certified object of
any kind attached to a ZEFOZ table (novelty sweep of 2026-08-12, shipped with
the release).  The claims here are application-first in that sense; every
mathematical tool is classical and credited above.

\section{The completeness question: a pre-registered kill}\label{sec:kill}

The completeness target was specified in advance: certify that the published
list is the whole list --- every stationary point of every $f_{ij}$, both
sites --- over the box $\|B\|_\infty\le100$~mT (the region the literature's own
grids explored; note that all twenty published points lie \emph{outside} this
box, at $191$~mT--$5.3$~T, so the completeness claim over the box would state
that $B=0$ is the only in-box stationary set).  The pre-registered kill
condition K2 read: at $\sim$20 laptop-hours of exclusion work, if the measured
statistics project more than $200$ laptop-hours to close the domain, kill the
completeness claim and downgrade to certified existence and curvature at the
known points.

The exclusion engine is the natural one: adaptive subdivision of the box; per
box $X=B_c+[-r,r]^3$, certified eigenvalue brackets at $B_c$ (Bauer--Fike
residual discs with component counting), per-level gap bounds over $X$ via the
Lipschitz bound $\sqrt3\,rD$, Davis--Kahan eigenvector enclosures over $X$,
certified Hellmann--Feynman gradient at $B_c$, a signed perturbation-sum
Hessian enclosure over $X$, and the first-order exclusion test $\lvert
g_k(B_c)\rvert>\sum_l \sup_X\lvert\partial_k\partial_l f\rvert\,r$ for some
$k$.  A box is closed when all 120 pairs are excluded; otherwise it splits.

The measured statistics are unambiguous.  Certified per-box cost is
$1.4$--$2.4$~s (all 120 pairs amortized over one diagonalization).  At box radius
$r=6.25$~mT and $r=0.78$~mT, \emph{zero} of the 120 pairs can be excluded at a
generic center ($B_c=(50,50,50)$~mT): the certified Hessian sup-bounds, inflated
by eigenvector drift $\propto rD/\gamma$, dominate the certified gradients.
First majority exclusions appear only at $r\approx0.1$~mT ($87/120$ at the
generic center; $68/120$ at $(6.25,6.25,6.25)$~mT), and $33\text{--}52$ pairs
survive even there.  Adaptive runs on sample chunks (six chunks of half-width
$0.39$~mT, both sites, stratified in radius) give
the following (all six exhausted their 340-s budgets with their
subdivision trees unfinished):
\begin{center}\small
\begin{tabular}{clccc}
\toprule
site & chunk center (mT) & boxes & closed boxes & surviving pairs \\
\midrule
1 & generic $(50,50,50)$ & 234 & 0 & 2 \\
1 & near-origin $(1.2,1.2,1.2)$ & 176 & 0 & 32 \\
1 & far $(87.5,62.5,37.5)$ & 177 & 153 & 0 \\
1 & near-axis $(75,5,5)$ & 136 & 10 & 3 \\
2 & generic $(50,50,50)$ & 251 & 218 & 0 \\
2 & near-origin $(1.2,1.2,1.2)$ & 191 & 0 & 19 \\
\bottomrule
\end{tabular}
\end{center}
Closed boxes are finest-level ($r=0.098$~mT) boxes with all 120 pairs
excluded; surviving pairs are those still active at the depth cap.  No
chunk finished, so the certified cost per unit volume strictly exceeds
each chunk's wall-clock/volume rate; the weakest measured rate is
$713$~s/mT$^3$.
Scaling that weakest rate to the half
domain ($4\times10^{6}$~mT$^3$ after the $B\to-B$ symmetry) gives a strict
lower bound of
more than $\mathbf{7.9\times10^{5}}$ per site ($1.58\times10^{6}$ for both sites)
laptop-hours --- before any refinement of the surviving pairs, which the
calibration shows needs two further depth levels for the flattest pairs at
generic fields and diverges for the near-origin web --- against the
pre-registered budget of $200$.  The lower bound alone exceeds the budget by
a factor of $8{,}000$, i.e.\ nearly four orders of magnitude.  \textbf{K2 fired; the completeness claim is dead at laptop scale},
and this note's certificate is the pre-registered downgrade.

The obstruction is structural and worth recording.  The zero-field spectrum of
each site contains near-degenerate level pairs --- certified splittings as
small as $0.004$~MHz (site~1) and $0.16$~MHz (site~2), with six site-1
spacings below $36$~MHz --- alongside spacings of $300$--$900$~MHz.  Every per-level certified bound in the exclusion engine ---
eigenvector drift, matrix-element enclosures, perturbation denominators ---
degrades as the ratio $rD/\gamma_n$, with $D\approx104$~MHz/mT the certified
field-Lipschitz constant and $\gamma_n$ the local gap of level $n$.  A level
inside a quasi-doublet of splitting $\gamma$ is processable only for
$r\lesssim\gamma/(4\sqrt3D)$, i.e.\ boxes of $10^{-5}$--$5\times10^{-2}$~mT
near the doublet-crossing set, which has positive two-dimensional measure in
the box.  Closing the domain therefore requires either cluster-projector
enclosures (certifying the two-dimensional doublet subspace as a block, with
block-resolved perturbation sums) or a different certification principle
altogether; with single-level machinery the cost diverges near the crossing
set.  We record this as the concrete engineering frontier, not as an
impossibility result.

\section{Errata in the reference}\label{sec:errata}

Three defects of \cite{Matsuura} were found and verified in the course of this
work.  We record them with version history; the first was corrected upstream
by the authors, and we credit that correction.

\begin{enumerate}
\item[(E1)] \emph{Site-1 quadrupole sign, v1 and v2.}  Appendix~B of v1
(2024-12-13) and v2 (2025-02-12) prints the site-1 $Q$ matrix with
$Q_{23}=Q_{32}=+15.5$~MHz.  With $+15.5$ the paper's own zero-field
frequencies (its Table~4) are irreproducible --- off by up to $68$~MHz ---
while with $-15.5$ they reproduce to $\le0.05$~MHz.  The paper's computations
therefore used $-15.5$ throughout; the printed table was wrong for eleven
months.  v3 (2025-11-14), which matches the version of record
\cite{MatsuuraPRB}, prints $-15.5$: corrected upstream, and any reader who
rationalized v1/v2 as printed inherited a wrong Hamiltonian.
\item[(E2)] \emph{Two tabulated frequencies inconsistent with their own
points, v3.}  In v3 Table~5, site~1, the pair $(6,7)$ is tabulated at
$745.8$~MHz and $(4,7)$ at $2216.2$~MHz.  The certified brackets at the
stationary points nearest the stated coordinates (Theorem~\ref{thm:points},
existence certified) are
$f_{67}=748.5431883\pm2\times10^{-10}$~MHz and
$f_{47}=2220.5387344\pm2\times10^{-10}$~MHz: discrepancies of $-2.74$ and
$-4.34$~MHz, roughly $50\times$ the rounding of the table.  No alternative
level pair at those fields is both stationary and at the tabulated frequency;
the other eighteen entries agree with our brackets to $\le0.07$~MHz (worst
$0.0634$~MHz, site~1 $(8,13)$).  The two values appear to be stale.
\item[(E3)] \emph{Sign pairing of the highlighted site-2 field, v3.}  The
high-precision bullet for the site-2 $(14,15)$ point prints
$B_{(D_1,D_2,b)}=(\mp378.9,\pm73.2,\mp502.3)$~mT alongside
$B_{(B,\theta,\phi)}=(633.52~\mathrm{mT},\mp37.5383^\circ,-10.9417^\circ)$.
The two are inconsistent: the stated angles (upper signs) give
$(-379.0,+73.3,+502.3)$~mT, and the certified stationary point
(Theorem~\ref{thm:points}) is at $(-378.99,+73.27,+502.35)$~mT.  The printed
Cartesian sign pairing should read $(\mp,\pm,\pm)$.
\end{enumerate}

We verified (E1) against the v1, v2, and v3 source files directly and (E2),
(E3) against certified brackets; the version of record \cite{MatsuuraPRB} was
not re-checked entry-by-entry against v3 beyond its Appendix~B.

\section{What is certified, what is enclosed, what is
trusted}\label{sec:trust}

\emph{Certified} (machine-checked from first principles by the standard-library
checker): every claim of Theorems~\ref{thm:zero} and \ref{thm:points} and
Corollary~\ref{cor:sig}, as inequalities among exact rationals: the symbolic
time-reversal identity, all inertia counts, all eigenvalue/transition brackets,
the gradient enclosures, the signed Hessian enclosures and their eigenvalue
brackets and signatures, and every Krawczyk contraction, including the
third-derivative bound of Appendix~\ref{app:third} re-derived from the
certified gaps.  \emph{Enclosed but pointing at published data:} the
identification of each certified point with a row of \cite[Table~5]{Matsuura}
is by proximity of the published rounded coordinates; the published rounding
($0.001$~T, $0.01^\circ$) is far coarser than our boxes.  \emph{Trusted:} the
Python \texttt{fractions}/\texttt{json}/\texttt{math} standard library; the
classical theorems cited (Sylvester, Davis--Kahan, Ostrowski, Krawczyk, and
the resolvent/Cauchy lemma proved in the appendix); and nothing else --- in
particular no floating point, no eigensolver, and no code shared with the
generator enters the verification.  The kill-condition statistics of
Section~\ref{sec:kill} are measurements of a private search engine and are
\emph{not} part of the certified surface; they are reported as run logs.

The computation, in one paragraph, per the program's protocol.  Candidate
fields and eigenvectors were produced by Newton iteration and 70-digit
diagonalization (mpmath), then frozen as dyadic rationals; certification is a
separate pass in exact rational interval arithmetic (outward rounding,
denominator cap $2^{200}$), $\sim4.6$~s per point; the 23-object certificate
(1.9~MB JSON: one time-reversal object, two zero-field objects, twenty
Krawczyk objects) verifies in $39$~s under \texttt{zefoz\_checker2.py}
(CPython~3.14, stdlib only, OS-independent), exit code 0; a six-item tamper
battery (corrupted inertia count, shifted Hessian bracket, shifted gradient
enclosure, flipped time-reversal sign, understated gradient norm, inflated
Krawczyk radius) is rejected item-by-item with exit code 1, and an untampered
control passes; the pilot certificate re-verifies (47~s, exit 0) and its 352
brackets contain an independent 60-digit recomputation.  The branch-and-bound
measurements ran as memory-capped queued jobs
(34~min wall-clock total, peak 38~MB per job against the 2500~MB cap); all artifacts, the checker, the run logs, and
the dated novelty sweep ship in the repository release; the search engine
itself stays private per the program's certificate-plus-checker model.

\section{Open questions}\label{sec:open}

\begin{question}
Cluster-certified completeness.  Can the box $\|B\|_\infty\le100$~mT be closed
by replacing per-level Davis--Kahan enclosures with certified two-dimensional
doublet-projector enclosures (block perturbation sums with inter-cluster gaps
only), and what is the measured cost?  The engine interface already isolates
the per-level bounds; this is the one concrete route our measurements leave
open at laptop scale.
\end{question}

\begin{question}
Existence certificates from double precision.  The Krawczyk test at clock
points needed candidates polished to $10^{-40}$; is there a certification
principle (exploiting the evenness of the spectrum in $B$, or the analyticity
of $f_{ij}$) whose stringency does not scale with the inverse Hessian, i.e.\
that certifies flat stationary points from double-precision data?
\end{question}

\begin{question}
The physical atlas.  The certified signatures (all twenty published
nonzero-field points saddles or maxima, none minima; six of the ten published
zero-field pairs saddles, the other four local minima) and the certified
zero-field curvatures are properties of the model \eqref{eq:H} with the
printed matrices.  Do measured second-order
shifts at the $(14,15)$ site-2 point distinguish the certified Hessian
spectrum $\{-1.652\times10^{-4},-8.949\times10^{-5},-1.447\times10^{-6}\}$
MHz/mT$^2$ from the scalar $\lvert S_2\rvert$ summaries in use?
\end{question}

\begin{question}
Which certified object does the experiment want next: brackets at the
$\lvert B\rvert\le3$~T points of the v3 search for site-independent field
orientations, certified \emph{optimality} of the $(10,11)$/$(14,15)$ choices
within the tabulated list, or completeness over the small ball
$\lvert B\rvert\le25$~mT actually swept by the published grids?
\end{question}

\appendix

\section{A third-derivative bound from a resolvent estimate}\label{app:third}

\begin{lemma}\label{lem:cauchy}
Let $H$ be Hermitian, $\lambda$ a simple eigenvalue with
$\gamma=\dist(\lambda,\operatorname{spec}H\setminus\{\lambda\})>0$, $V$
Hermitian with $\|V\|\le D$.  For complex $t$ with $\lvert t\rvert<\rho:=
\gamma/(4D)$ the matrix $H+tV$ has exactly one eigenvalue $\lambda(t)$ in the
open disc of radius $\gamma/2$ about $\lambda$; $\lambda(t)$ is analytic, and
$\bigl|\lambda'''(0)\bigr|\le 192\,D^3/\gamma^2$.
\end{lemma}

\begin{proof}
This is the standard resolvent perturbation argument (cf.\ \cite{Kato}) with
constants tracked.  For $z$ on the circle $\Gamma$ of radius $\gamma/2$ about
$\lambda$,
$\|(H-z)^{-1}\|=\dist(z,\operatorname{spec}H)^{-1}\le2/\gamma$, so for
$\lvert t\rvert<\rho$ the Neumann series of
$(H+tV-z)^{-1}=(H-z)^{-1}\sum_{k\ge0}\bigl(-tV(H-z)^{-1}\bigr)^k$ converges
($\lVert tV(H-z)^{-1}\rVert\le\rho D\cdot2/\gamma=1/2$).  Hence $\Gamma$ lies
in the resolvent set of $H+tV$ for all such $t$; the Riesz projection
$P(t)=-\frac1{2\pi i}\oint_\Gamma(H+tV-z)^{-1}dz$ is analytic in $t$ and of
constant rank $1$, and $\lambda(t)=\tr\bigl[(H+tV)P(t)\bigr]$ is the unique
eigenvalue inside $\Gamma$, analytic in $t$, with
$\lvert\lambda(t)-\lambda\rvert\le\gamma/2$.  Writing
$\lambda(t)-\lambda=\sum_{k\ge1}a_kt^k$ on $\lvert t\rvert<\rho$ and applying
the Cauchy estimates on circles of radius $\rho'\uparrow\rho$,
$\lvert a_3\rvert\le(\gamma/2)/\rho^3$, so
$\lvert\lambda'''(0)\rvert=6\lvert a_3\rvert\le3\gamma\,(4D/\gamma)^3
=192\,D^3/\gamma^2$.
\end{proof}

Applied at every real point $B$ of a box $X$ with $V=\partial_e H=\sum_k e_k
Z_k$ ($\lvert e\rvert=1$, so $\|V\|\le D:=(\sum_k\|Z_k\|_2^2)^{1/2}$) and with
$\gamma$ replaced by certified box-uniform gap lower bounds
$\tilde\gamma_i,\tilde\gamma_j$, the lemma gives, for the cubic form
$g_B(e)=\partial_e^3 f_{ij}(B)$,
$\sup_{\lvert e\rvert=1}\lvert g_B(e)\rvert\le
M:=192\,D^3(\tilde\gamma_i^{-2}+\tilde\gamma_j^{-2})$.  Polarization of the
symmetric trilinear form $T=D^3f_{ij}(B)$,
$T[a,b,c]=\frac1{48}\sum_{\varepsilon\in\{\pm1\}^3}\varepsilon_1\varepsilon_2
\varepsilon_3\,g_B(\varepsilon_1a+\varepsilon_2b+\varepsilon_3c)$, with
$\lVert\varepsilon_1a+\varepsilon_2b+\varepsilon_3c\rVert\le3$ for unit
$a,b,c$, gives $\lvert T[a,b,c]\rvert\le\frac{8\cdot27}{48}M=4.5\,M$; the
certificate uses the laxer constant $9M$.  Integrating along the segment from
$B^*$ to $B\in X$ bounds each Hessian entry's variation by
$9M\lVert B-B^*\rVert\le9\sqrt3\,Mr$, which is the entrywise inflation used in
the Krawczyk operator, and which the checker recomputes from the certified
$D$ and gaps.

\begin{thebibliography}{99}

\bibitem{Matsuura} K.~Matsuura, S.~Yasui, R.~Kaji, H.~Sasakura, T.~Tawara,
S.~Adachi, \emph{Exploration of optimal hyperfine transitions for spin-wave
storage in $^{167}$Er$^{3+}$:Y$_2$SiO$_5$}, arXiv:2412.10126; v1 2024-12-13,
v2 2025-02-12, v3 2025-11-14.  Version citations in the text refer to the
arXiv versions.

\bibitem{MatsuuraPRB} K.~Matsuura et al., Phys.\ Rev.\ B \textbf{113}, 085421
(2026).  Version of record of \cite{Matsuura}.

\bibitem{Wang2023} S.-J.~Wang, Y.-H.~Chen, J.~J.~Longdell, X.~Zhang,
J.~Lumin.\ \textbf{262}, 119935 (2023).  Source of the $A$, $Q$, $g$ matrices
reproduced in \cite[App.~B]{Matsuura}.

\bibitem{Fraval2004} E.~Fraval, M.~J.~Sellars, J.~J.~Longdell, \emph{Method of
extending hyperfine coherence times in ${\rm Pr}^{3+}$:${\rm Y_2SiO_5}$},
Phys.\ Rev.\ Lett.\ \textbf{92}, 077601 (2004).

\bibitem{Zhong2015} M.~Zhong et al., \emph{Optically addressable nuclear spins
in a solid with a six-hour coherence time}, Nature \textbf{517}, 177 (2015).

\bibitem{Ortu2601} Z.~Wang et al., \emph{Quadrupole coupling and ZEFOZ points
of $^{167}$Er in CaWO$_4$}, arXiv:2601.16362 (2026).

\bibitem{Business2606} \emph{On-demand coherent mapping of telecom optical
states onto erbium hyperfine spins}, arXiv:2606.15009 (2026).

\bibitem{DavisKahan} C.~Davis, W.~M.~Kahan, \emph{The rotation of eigenvectors
by a perturbation. III}, SIAM J.\ Numer.\ Anal.\ \textbf{7}, 1--46 (1970).

\bibitem{Krawczyk} R.~Krawczyk, \emph{Newton-Algorithmen zur Bestimmung von
Nullstellen mit Fehlerschranken}, Computing \textbf{4}, 187--201 (1969).

\bibitem{Neumaier} A.~Neumaier, \emph{Interval Methods for Systems of
Equations}, Cambridge University Press (1990).

\bibitem{HornJohnson} R.~A.~Horn, C.~R.~Johnson, \emph{Matrix Analysis}, 2nd
ed., Cambridge University Press (2013).  Ostrowski's theorem is Thm.~4.5.9.

\bibitem{Kato} T.~Kato, \emph{Perturbation Theory for Linear Operators},
Springer (1980).

\bibitem{Rump} S.~M.~Rump, \emph{Verification methods: rigorous results using
floating-point arithmetic}, Acta Numerica \textbf{19}, 287--449 (2010).

\bibitem{symveig} \emph{symveig: verified eigenvalue enclosures for
symmetry-decomposed Hermitian matrices}, arXiv:2606.16217 (2026).

\end{thebibliography}

\end{document}
